%====================================================================
% Rapport DAAR – Projet 2 : Moteur de recherche de bibliothèque
% Université Sorbonne – M1 Informatique
% Auteurs : Lina Tireche & Nouha Kara
%====================================================================
\documentclass[12pt,a4paper]{article}

%---- Encodage & langue ----
\usepackage[utf8]{inputenc}
\usepackage[T1]{fontenc}
\usepackage[french]{babel}
\usepackage{lmodern}

%---- Mise en page ----
\usepackage[top=2.5cm, bottom=2.5cm, left=2.5cm, right=2.5cm]{geometry}
\usepackage{microtype}
\usepackage{parskip}
\usepackage{setspace}
\onehalfspacing

%---- En-têtes / pieds de page ----
\usepackage{fancyhdr}
\pagestyle{fancy}
\fancyhf{}
\fancyhead[L]{\small DAAR – Projet 2}
\fancyhead[R]{\small Tireche \& Kara}
\fancyfoot[C]{\thepage}
\renewcommand{\headrulewidth}{0.4pt}

%---- Maths (minimaliste) ----
\usepackage{amsmath}

%---- Tableaux ----
\usepackage{booktabs}
\usepackage{tabularx}
\usepackage{array}

%---- Listings (code) ----
\usepackage{listings}
\lstset{
  basicstyle=\ttfamily\small,
  frame=single,
  breaklines=true,
  numbers=left,
  numberstyle=\tiny,
  stepnumber=1,
  tabsize=4,
  showstringspaces=false,
  keywordstyle=\bfseries,
  commentstyle=\itshape,
}
\lstdefinestyle{python}{language=Python}
\lstdefinestyle{bash}{language=bash}
\lstdefinestyle{yaml}{language=Python,
  morekeywords={true,false,null,on,off,yes,no}}

%---- TikZ ----
\usepackage{tikz}
\usetikzlibrary{shapes.geometric, arrows.meta, positioning}

%---- Liens (noir, sans couleur) ----
\usepackage[hidelinks]{hyperref}

%---- Utilitaires ----
\usepackage{graphicx}
\usepackage{enumitem}

%====================================================================
\begin{document}

%---- Page de titre ----
\begin{titlepage}
  \centering
  \vspace*{2cm}
  {\Large \textbf{Sorbonne Université}\\[0.3em]
   \large UFR d'Informatique – Master 1 Informatique}\\[1cm]

  \rule{\linewidth}{0.8pt}\\[0.6em]
  {\LARGE \textbf{Projet 2 – Algorithmique des Réseaux et Données}}\\[0.4em]
  {\large \textit{Moteur de recherche de bibliothèque textuelle}}\\[0.3em]
  \rule{\linewidth}{0.8pt}\\[2cm]

  {\large \textbf{Auteurs :}\\[0.4em]
   Lina \textsc{Tireche}\\[0.2em]
   Nouha \textsc{Kara}}\\[2cm]

  {\large \textbf{Encadrant :}\\[0.4em]
   Binh-Minh \textsc{Bui-Xuan}}\\[2cm]

  {\large Année universitaire 2025–2026}\\[1cm]

  \vfill
  {\small Dépôt Git : \texttt{https://github.com/linatireche/daar-projet2}}
\end{titlepage}

%---- Résumé ----
\begin{abstract}
Ce rapport présente la réalisation complète d'un moteur de recherche textuelle
appliqué à un corpus de livres issus du Projet Gutenberg.
Nous avons construit une application web full-stack composée d'un backend
Python (FastAPI), d'un frontend React et d'une base de données PostgreSQL.
L'application permet de rechercher des livres par mots-clés, de recommander
des ouvrages similaires et de classer les livres selon leur importance dans
un graphe de proximité thématique. Le tout est containerisé avec Docker,
déployable sur Kubernetes, et soumis à une intégration continue via GitHub Actions.
\end{abstract}

\tableofcontents
\newpage

%====================================================================
\section{Introduction}
%====================================================================

L'objectif de ce projet était de concevoir et développer, de bout en bout,
un moteur de recherche capable d'interroger une bibliothèque numérique de grande
taille. Nous avons téléchargé automatiquement plus de 1\,600 livres depuis le
Projet Gutenberg, puis nous les avons indexés et analysés à l'aide de plusieurs
algorithmes étudiés en cours.

Dès le départ, nous avons fait le choix de construire une véritable application
web, et pas seulement une série de scripts. Cela signifiait concevoir une
architecture cohérente, gérer une base de données, exposer une API REST,
développer une interface utilisateur et enfin déployer le tout dans des conteneurs.
Ce choix nous a demandé un effort supplémentaire, mais il nous a permis de
comprendre comment les algorithmes s'intègrent dans un système réel.

Le rapport est organisé de la façon suivante : nous commençons par décrire
l'architecture générale du projet, puis nous détaillons chaque fonctionnalité
implémentée — la recherche de texte, le graphe de similarité, l'analyse du graphe —
et nous terminons par les aspects déploiement, tests et interface utilisateur.

%====================================================================
\section{Architecture du projet}
%====================================================================

\subsection{Vue d'ensemble}

Le projet est structuré en trois grandes couches qui communiquent entre elles :
le frontend, le backend et la base de données. Cette séparation nous a permis
de travailler en parallèle sur les différentes parties et de remplacer ou
modifier un composant sans tout casser.

\begin{figure}[h]
\centering
\begin{tikzpicture}[
  node distance=1.6cm and 2.8cm,
  box/.style={rectangle, draw, thick, minimum width=3.2cm,
              minimum height=1.1cm, align=center, font=\small},
  arr/.style={-{Stealth}, thick}
]
  \node[box] (fe)  {Frontend\\React (port 80)};
  \node[box, right=of fe] (be) {Backend\\FastAPI (port 8000)};
  \node[box, right=of be] (db) {PostgreSQL\\(port 5432)};

  \draw[arr] (fe) -- (be) node[midway, above, font=\scriptsize]{REST/HTTP};
  \draw[arr] (be) -- (db) node[midway, above, font=\scriptsize]{SQL};
  \draw[arr, dashed] (be.north) -- ++(0,1)
        node[above, font=\scriptsize]{API Gutenberg};
\end{tikzpicture}
\caption{Architecture trois-tiers de l'application.}
\label{fig:archi}
\end{figure}

\subsection{Le backend — FastAPI}

Le backend est écrit en Python avec le framework FastAPI. C'est lui qui
contient toute la logique métier : il reçoit les requêtes du frontend,
exécute les algorithmes et renvoie les résultats au format JSON.

Nous avons organisé le code backend en plusieurs modules bien séparés.
Le dossier \texttt{algorithms/} contient les implémentations pures des
algorithmes (KMP, index inversé, Jaccard, PageRank, etc.).
Le dossier \texttt{api/} contient les routes HTTP.
Le dossier \texttt{services/} fait le lien entre les deux : il appelle
les algorithmes et les fonctions de base de données, et formate les réponses.
Enfin, le dossier \texttt{models/} définit les structures de données
(modèles SQLAlchemy et schémas Pydantic).

\begin{lstlisting}[caption={Arborescence du backend.}, label=lst:tree]
backend/
|-- app/
|   |-- algorithms/
|   |   |-- kmp.py
|   |   |-- inverted_index.py
|   |   |-- jaccard_graph.py
|   |   `-- centrality.py
|   |-- api/
|   |   |-- books.py
|   |   |-- search.py
|   |   `-- stats.py
|   |-- models/
|   |   `-- book.py
|   |-- services/
|   |   `-- book_service.py
|   `-- database.py
`-- tests/
    `-- test_algorithms.py
\end{lstlisting}

\subsection{Le frontend — React}

Le frontend est une application React développée en TypeScript avec Vite
comme outil de build. Il communique avec le backend uniquement via des
appels \texttt{fetch} vers l'API REST — il n'y a aucun état partagé côté serveur.
L'application est servie par Nginx une fois compilée, ce qui la rend très légère
et rapide.

\subsection{La base de données — PostgreSQL}

Nous utilisons PostgreSQL 16 pour stocker les livres, leur contenu, leurs
métadonnées (auteur, titre, langue, nombre de mots) ainsi que les arêtes
du graphe de similarité et les scores de centralité. Le choix de PostgreSQL
nous a semblé naturel : c'est une base robuste, très bien supportée par
SQLAlchemy (l'ORM Python que nous utilisons), et parfaitement intégrée
dans l'écosystème Docker.

Toute la communication entre le backend et la base de données passe par
SQLAlchemy, ce qui nous évite d'écrire du SQL brut et rend le code plus
lisible et plus maintenable.

%====================================================================
\section{Collecte et préparation des données}
%====================================================================

\subsection{Téléchargement des livres}

Les livres proviennent du Projet Gutenberg, une bibliothèque numérique qui
met à disposition des milliers d'œuvres du domaine public. Nous avons utilisé
leur API pour télécharger les textes automatiquement au démarrage de l'application.

Le service de collecte récupère les livres par batch, en gérant les limites
de débit de l'API (rate limiting). Chaque livre est enregistré en base avec
son identifiant Gutenberg, son titre, son auteur, sa langue et son contenu
complet en texte brut.

\subsection{Prétraitement du texte}

Avant d'indexer les livres, nous appliquons quelques étapes de nettoyage :
suppression des en-têtes et pieds de page standard de Gutenberg,
normalisation des espaces et de la ponctuation, conversion en minuscules.
Pour la construction du graphe de similarité, nous filtrons également les
mots courts (moins de 4 caractères) et les mots très fréquents qui ne sont
pas discriminants (stopwords).

\subsection{Résultats de la collecte}

Au final, notre corpus contient \textbf{1\,665 livres}, pour un total de
près de \textbf{18,8 millions de mots}. La longueur moyenne d'un livre est
d'environ 11\,300 mots. Ces chiffres sont affichés en temps réel dans le
tableau de bord de l'application.

%====================================================================
\section{Recherche dans les textes}
%====================================================================

\subsection{L'algorithme de Knuth-Morris-Pratt (KMP)}

\subsubsection{Ce que fait KMP}

Lorsqu'un utilisateur saisit un mot ou une expression dans la barre de
recherche et choisit le mode « recherche exacte », notre application utilise
l'algorithme KMP pour parcourir le contenu de chaque livre et trouver toutes
les occurrences de ce motif.

L'intérêt de KMP par rapport à une recherche naïve est qu'il ne reprend jamais
depuis le début en cas d'échec partiel. Il mémorise, grâce à une table de
préfixes calculée à l'avance, jusqu'où il peut \og sauter\fg{} dans le motif
sans rater une occurrence. Concrètement, cela lui permet de parcourir le texte
en un seul passage, sans jamais reculer.

\subsubsection{Comment nous l'avons implémenté}

Nous avons implémenté KMP from scratch en Python dans le fichier
\texttt{algorithms/kmp.py}. La fonction prend en entrée le texte d'un livre
et le motif recherché, et renvoie la liste des positions (indices de caractères)
où le motif apparaît.

\begin{lstlisting}[style=python, caption={Implémentation de KMP.}, label=lst:kmp]
def build_failure(pattern: str) -> list[int]:
    """Construit la table de prefixes (table pi) du motif."""
    m = len(pattern)
    pi = [0] * m
    k = 0
    for i in range(1, m):
        while k > 0 and pattern[k] != pattern[i]:
            k = pi[k - 1]
        if pattern[k] == pattern[i]:
            k += 1
        pi[i] = k
    return pi

def kmp_search(text: str, pattern: str) -> list[int]:
    """Renvoie les positions de toutes les occurrences de pattern dans text."""
    if not pattern:
        return []
    pi = build_failure(pattern)
    matches, q = [], 0
    for i, c in enumerate(text):
        while q > 0 and pattern[q] != c:
            q = pi[q - 1]
        if pattern[q] == c:
            q += 1
        if q == len(pattern):
            matches.append(i - len(pattern) + 1)
            q = pi[q - 1]
    return matches
\end{lstlisting}

En plus des positions brutes, le service associé extrait un extrait de texte
(snippet) autour de chaque occurrence, que l'interface affiche pour donner
du contexte à l'utilisateur.

\subsection{L'index inversé avec TF-IDF}

\subsubsection{Ce que fait l'index inversé}

La recherche par KMP est précise mais lente sur un grand corpus : elle doit
lire chaque livre en entier pour chaque requête. L'index inversé résout ce
problème en préparant à l'avance une structure qui associe chaque mot à la
liste des livres qui le contiennent.

Quand un utilisateur cherche \og adventure\fg{}, au lieu de lire 1\,665 livres,
on consulte directement l'entrée de l'index pour ce mot, et on obtient
instantanément la liste des livres concernés. C'est le même principe que
l'index à la fin d'un manuel scolaire.

Nous avons enrichi cet index avec un score TF-IDF, qui mesure à quel point
un mot est important dans un livre donné. Un mot très fréquent dans un livre
mais rare dans le reste du corpus aura un score élevé — c'est un bon indicateur
que le livre traite vraiment de ce sujet.

\subsubsection{Comment nous l'avons implémenté}

L'index est construit une seule fois au démarrage du serveur, en parcourant
tous les livres. Il est ensuite conservé en mémoire sous forme d'un dictionnaire
Python. Cette construction prend quelques secondes, mais ensuite chaque
recherche est quasiment instantanée.

\subsubsection{Comparaison des performances}

Nous avons mesuré les temps de réponse moyens sur notre corpus :

\begin{table}[h]
\centering
\caption{Comparaison des méthodes de recherche.}
\label{tab:bench}
\begin{tabular}{lrp{6cm}}
\toprule
\textbf{Méthode} & \textbf{Temps moyen} & \textbf{Avantage principal} \\
\midrule
Index inversé (TF-IDF) & 6 ms  & Résultats triés par pertinence, très rapide \\
Expressions régulières  & 8 ms  & Syntaxe flexible (wildcards, etc.) \\
KMP                     & 26 ms & Recherche exacte, toutes occurrences garanties \\
\bottomrule
\end{tabular}
\end{table}

L'index inversé est 4 fois plus rapide que KMP car il ne lit jamais les
livres lors de la requête. KMP est plus lent mais il est le seul à trouver
précisément toutes les occurrences d'une expression dans le texte brut.

%====================================================================
\section{Graphe de similarité — Jaccard}
%====================================================================

\subsection{L'idée générale}

Une des fonctionnalités les plus intéressantes de l'application est la
recommandation : quand un utilisateur consulte un livre, l'application lui
propose d'autres livres similaires. Pour cela, nous avons construit un graphe
dans lequel chaque nœud est un livre et chaque arête relie deux livres
thématiquement proches.

La mesure de proximité que nous utilisons est la similarité de Jaccard.
Elle consiste à comparer les vocabulaires de deux livres : deux livres qui
partagent beaucoup de mots en commun (relativement à l'union de leurs
vocabulaires) sont considérés comme proches.

\subsection{Comment nous avons construit le graphe}

Pour chaque livre, nous commençons par extraire son vocabulaire significatif :
on tokenise le texte (on découpe en mots), on filtre les mots trop courts
et les stopwords, et on garde l'ensemble des mots distincts.

Ensuite, pour chaque paire de livres, on calcule leur similarité de Jaccard.
Si cette similarité dépasse un seuil de 0,12 (que nous avons fixé après
expérimentation), on ajoute une arête entre ces deux livres dans le graphe.
Pour éviter d'avoir des graphes trop denses, nous limitons à 6 le nombre de
voisins par livre (on garde seulement les 6 livres les plus similaires).

Cette étape de construction est la plus coûteuse en temps de calcul, car
elle compare toutes les paires possibles. Avec 1\,665 livres, cela représente
plus d'un million de comparaisons. C'est pourquoi nous l'exécutons une seule
fois au démarrage et nous persistons les résultats en base de données.

\subsection{Résultats}

Le graphe final contient \textbf{9\,628 arêtes}. Ces connexions reflètent
bien les proximités thématiques : par exemple, les romans d'aventure sont
reliés entre eux, les œuvres scientifiques se regroupent, et les textes
philosophiques forment leur propre cluster.

\begin{lstlisting}[style=python, caption={Calcul de la similarité de Jaccard.}, label=lst:jaccard]
def jaccard(set_a: set, set_b: set) -> float:
    """Retourne la similarite de Jaccard entre deux ensembles de mots."""
    if not set_a and not set_b:
        return 1.0
    intersection = len(set_a & set_b)
    union = len(set_a | set_b)
    return intersection / union if union else 0.0
\end{lstlisting}

%====================================================================
\section{Analyse du graphe — Centralités}
%====================================================================

Une fois le graphe construit, nous l'analysons pour calculer l'importance
relative de chaque livre. Nous avons implémenté trois mesures de centralité
classiques, chacune donnant une lecture différente du graphe.

\subsection{PageRank}

Le PageRank est l'algorithme qu'utilisait Google à l'origine pour classer
les pages web. Appliqué à notre graphe de livres, il mesure l'importance
d'un livre en tenant compte non seulement du nombre de ses voisins, mais
aussi de l'importance de ces voisins.

Concrètement, un livre a un PageRank élevé si beaucoup de livres importants
lui sont reliés. C'est une mesure récursive : on part d'une distribution
uniforme de scores, et on les propage de proche en proche jusqu'à convergence.
En pratique, sur notre graphe, l'algorithme converge en moins de 35 itérations.

Nous avons implémenté PageRank from scratch en Python, avec un facteur
d'amortissement de 0,85 (valeur standard) et un critère d'arrêt basé sur
la variation des scores entre deux itérations consécutives.

Les livres avec le PageRank le plus élevé sont en quelque sorte les
\og classiques\fg{} du corpus : des œuvres très centrales, recommandées
par de nombreux autres livres eux-mêmes importants.

\subsection{Closeness centrality}

La closeness mesure à quel point un livre est \og proche\fg{} de tous les autres
dans le graphe, en termes de nombre d'arêtes à traverser. Un livre avec une
closeness élevée peut \og atteindre\fg{} n'importe quel autre livre en peu d'étapes.

Nous avons implémenté ce calcul par des parcours en largeur (BFS) depuis
chaque nœud du graphe. Comme notre graphe peut avoir plusieurs composantes
connexes (des livres sans voisins), nous utilisons la correction de
Wasserman-Faust qui ramène les scores à une échelle comparable même pour
les nœuds isolés.

Les livres à forte closeness sont ceux qui ont des thèmes polyvalents,
qui leur permettent d'être reliés à des ouvrages de genres variés.

\subsection{Betweenness centrality}

La betweenness mesure à quel point un livre sert de \og pont\fg{} dans le graphe.
Plus précisément, elle compte la proportion de plus courts chemins entre
toutes les paires de livres qui passent par ce livre.

Un livre avec une betweenness élevée joue un rôle structurant dans le graphe :
sans lui, certains clusters de livres se trouveraient déconnectés les uns des
autres. Typiquement, ce sont des livres qui traitent de plusieurs thèmes à la fois.

Nous avons utilisé l'algorithme de Brandes pour calculer la betweenness de
tous les nœuds en un seul passage, ce qui est bien plus efficace que de
calculer naïvement les plus courts chemins entre toutes les paires.

\subsection{Exposition via l'API}

Les trois scores sont calculés au démarrage du serveur et stockés en base
de données avec le livre correspondant. Ils sont exposés via la route
\texttt{GET /api/books/ranked?metric=pagerank} (ou \texttt{closeness},
ou \texttt{betweenness}), que le frontend utilise pour afficher la vue
\og Classement\fg{}.

%====================================================================
\section{Déploiement}
%====================================================================

\subsection{Containerisation avec Docker}

L'une de nos contraintes était que l'application doit pouvoir tourner
n'importe où, sans installation manuelle de dépendances. Nous avons donc
containerisé les trois composants (frontend, backend, base de données)
avec Docker.

Chaque service est décrit dans un fichier \texttt{Dockerfile} :
\begin{itemize}
  \item Le Dockerfile du backend part d'une image Python 3.11 légère,
        installe les dépendances listées dans \texttt{requirements.txt} et
        lance le serveur Uvicorn.
  \item Le Dockerfile du frontend utilise une construction multi-étape :
        d'abord Node 20 pour compiler l'application React avec Vite,
        ensuite Nginx Alpine pour servir les fichiers statiques générés.
        L'image finale pèse à peine quelques mégaoctets.
  \item La base de données utilise directement l'image officielle
        PostgreSQL 16, sans personnalisation.
\end{itemize}

Le fichier \texttt{docker-compose.yml} orchestre les trois services,
définit les variables d'environnement, les ports exposés et les dépendances
de démarrage (le backend attend que la base soit prête avant de lancer).

\begin{lstlisting}[style=yaml, caption={Extrait de docker-compose.yml.}, label=lst:compose]
services:
  db:
    image: postgres:16-alpine
    environment:
      POSTGRES_DB: daar_library
      POSTGRES_USER: daar_user
      POSTGRES_PASSWORD: daar_password
    volumes:
      - pg_data:/var/lib/postgresql/data
    healthcheck:
      test: ["CMD-SHELL", "pg_isready -U daar_user"]
      interval: 5s
      retries: 5

  backend:
    build: ./backend
    ports:
      - "8000:8000"
    depends_on:
      db:
        condition: service_healthy
    env_file: .env

  frontend:
    build: ./frontend
    ports:
      - "80:80"
    depends_on:
      - backend

volumes:
  pg_data:
\end{lstlisting}

Avec cette configuration, lancer toute l'application se résume à une seule
commande : \texttt{docker compose up --build}.

\subsection{Déploiement sur Kubernetes}

En plus de Docker Compose (pour le développement local), nous avons préparé
des manifests Kubernetes pour un déploiement en production.

Le dossier \texttt{k8s/} contient les fichiers YAML suivants :
\begin{itemize}
  \item Un \texttt{Deployment} et un \texttt{Service} pour chacun des trois
        composants (frontend, backend, PostgreSQL).
  \item Un \texttt{PersistentVolumeClaim} pour la persistance des données
        PostgreSQL, même si le pod redémarre.
  \item Un \texttt{Ingress} NGINX qui route le trafic HTTP entrant : les
        requêtes vers \texttt{/api/} sont redirigées vers le backend,
        tout le reste va vers le frontend.
  \item Un \texttt{ConfigMap} pour les variables de configuration et un
        \texttt{Secret} pour les informations sensibles (mot de passe de
        la base de données).
\end{itemize}

Nous avons également écrit un script \texttt{k8s-deploy.sh} qui automatise
l'application de tous ces manifests dans le bon ordre.

\subsection{Pipeline CI/CD — GitHub Actions}

Chaque fois que nous poussons du code sur la branche principale ou que nous
ouvrons une pull request, un workflow GitHub Actions se déclenche
automatiquement. Ce pipeline réalise les étapes suivantes :

\begin{enumerate}
  \item Il configure un environnement Python 3.11 et installe toutes les
        dépendances du backend.
  \item Il démarre une instance PostgreSQL éphémère (fournie comme service
        par GitHub Actions) pour que les tests puissent s'appuyer sur une
        vraie base de données.
  \item Il exécute la suite complète de tests pytest.
  \item Il vérifie que le frontend se compile sans erreur avec
        \texttt{npm run build}.
\end{enumerate}

Si l'une de ces étapes échoue, le push est signalé comme défaillant sur GitHub
et nous sommes immédiatement alertées. Cela nous a permis de détecter plusieurs
régressions pendant le développement.

%====================================================================
\section{Interface utilisateur}
%====================================================================

L'interface React est organisée autour de quatre vues principales, accessibles
depuis une barre de navigation commune.

\subsection{La recherche}

C'est la page d'accueil de l'application. L'utilisateur saisit un ou plusieurs
mots dans une barre de recherche et choisit la méthode souhaitée parmi trois
options : index inversé (TF-IDF), KMP ou expressions régulières.

Les résultats s'affichent sous forme de cartes, chacune montrant le titre,
l'auteur, un extrait du texte autour de la correspondance, et le score de
pertinence (pour la recherche par index). Un clic sur une carte ouvre la vue
détail du livre.

\subsection{Le détail d'un livre}

Cette page affiche les métadonnées complètes du livre (titre, auteur, langue,
nombre de mots, date de publication) ainsi que le début du texte.

En bas de page, une section \textit{Livres similaires} présente jusqu'à six
recommandations calculées par la similarité de Jaccard. Chaque recommandation
affiche le score de similarité et permet de naviguer vers le livre suggéré.

\subsection{Le classement}

Cette vue affiche la liste de tous les livres triés par l'une des trois mesures
de centralité, sélectionnable via un menu déroulant : PageRank, closeness ou
betweenness. Le score associé est affiché pour chaque livre.

\subsection{Les statistiques}

Le tableau de bord affiche des indicateurs globaux sur le corpus et le graphe :
nombre total de livres, nombre total de mots, taille du vocabulaire, nombre
d'arêtes dans le graphe de Jaccard, densité du graphe, et les dix livres les
mieux classés par PageRank. Ces données sont utiles pour vérifier que l'indexation
s'est bien déroulée.

%====================================================================
\section{Tests}
%====================================================================

Nous avons mis en place une suite de 7 tests automatisés avec \texttt{pytest}.
Ces tests couvrent à la fois les algorithmes purs et les routes de l'API.

\begin{table}[h]
\centering
\caption{Suite de tests automatisés.}
\label{tab:tests}
\begin{tabularx}{\linewidth}{lXl}
\toprule
\textbf{Test} & \textbf{Ce qui est vérifié} & \textbf{Résultat} \\
\midrule
\texttt{test\_kmp\_found}
  & KMP trouve toutes les occurrences d'un motif présent dans le texte
  & \checkmark \\[4pt]
\texttt{test\_kmp\_not\_found}
  & KMP renvoie une liste vide si le motif est absent
  & \checkmark \\[4pt]
\texttt{test\_index\_build}
  & L'index inversé indexe correctement les termes d'un mini-corpus
  & \checkmark \\[4pt]
\texttt{test\_index\_query}
  & Une requête TF-IDF renvoie le bon document en tête
  & \checkmark \\[4pt]
\texttt{test\_jaccard}
  & La similarité entre deux ensembles connus est correcte
  & \checkmark \\[4pt]
\texttt{test\_pagerank}
  & La somme des scores PageRank est bien égale à 1 (conservation)
  & \checkmark \\[4pt]
\texttt{test\_api\_books}
  & La route \texttt{GET /api/books} retourne un code 200
  & \checkmark \\
\bottomrule
\end{tabularx}
\end{table}

Les tests unitaires (KMP, index, Jaccard, PageRank) n'ont pas besoin de base
de données et s'exécutent très rapidement. Le test d'intégration
(\texttt{test\_api\_books}) utilise une base de données de test éphémère,
automatiquement créée et détruite par le fixture pytest correspondant.

Tous les tests passent en continu depuis la mise en place du pipeline CI/CD.

%====================================================================
\section{Difficultés rencontrées et solutions}
%====================================================================

Le développement de ce projet ne s'est pas fait sans obstacles. Voici les
principaux problèmes que nous avons rencontrés et comment nous les avons résolus.

\textbf{La connexion à la base de données dans Docker.}
Au début, le backend utilisait \texttt{localhost} comme adresse de la base,
ce qui fonctionnait en local mais pas dans Docker Compose, où la base est
accessible via son nom de service (\texttt{db}). Nous avons centralisé la
configuration dans un fichier \texttt{.env} et harmonisé toutes les références
à l'URL de connexion.

\textbf{L'initialisation de la session SQLAlchemy.}
Nous avons eu un problème avec l'initialisation de la session SQLAlchemy :
la variable \texttt{SessionLocal} était parfois \texttt{None} lors du premier
appel. Nous avons résolu ce problème en refactorisant la gestion de l'état de
l'application avec un singleton \texttt{AppState}, initialisé une seule fois
au démarrage du serveur.

\textbf{La migration vers Pydantic v2.}
La version 2 de Pydantic (utilisée par les versions récentes de FastAPI) a
introduit des changements incompatibles avec v1. En particulier, la méthode
\texttt{from\_orm()} a été remplacée par \texttt{model\_validate()}.
Nous avons mis à jour tous les appels correspondants dans \texttt{book\_service.py}
et \texttt{stats.py}.

\textbf{La performance de la construction du graphe.}
Avec 1\,665 livres, la comparaison de toutes les paires prend un temps non
négligeable. Nous avons optimisé en pré-calculant les ensembles de mots
une seule fois par livre et en parallélisant les comparaisons. Les résultats
sont ensuite persistés en base pour ne pas refaire le calcul à chaque démarrage.

%====================================================================
\section{Conclusion}
%====================================================================

Ce projet nous a donné l'opportunité de mettre en pratique, dans un contexte
réaliste, plusieurs algorithmes fondamentaux vus en cours : la recherche de
chaînes avec KMP, l'indexation avec TF-IDF, la mesure de similarité avec
Jaccard, et l'analyse de graphes avec PageRank, closeness et betweenness.

Au-delà des algorithmes eux-mêmes, nous avons dû concevoir une architecture
logicielle complète, gérer une base de données, construire une API REST,
développer une interface utilisateur, et déployer l'ensemble dans des conteneurs.
Ce côté \og ingénierie logicielle\fg{} était aussi enrichissant que le côté algorithmique.

Les résultats sont concluants : l'application tourne de manière stable,
les recherches répondent en quelques millisecondes, les recommandations sont
cohérentes, et le pipeline CI/CD garantit que les tests passent à chaque
modification du code.

Des améliorations resteraient envisageables : utiliser MinHash pour accélérer
le calcul de Jaccard à grande échelle, ajouter un moteur de suggestion
orthographique basé sur la distance de Levenshtein, ou mettre en cache les
recherches fréquentes avec Redis.

%====================================================================
\begin{thebibliography}{9}

\bibitem{kmp1977}
  D.~E. Knuth, J.~H. Morris et V.~R. Pratt,
  \textit{Fast Pattern Matching in Strings},
  SIAM Journal on Computing, vol.~6, n\textsuperscript{o}~2,
  pp.~323--350, 1977.

\bibitem{jaccard1912}
  P.~Jaccard,
  \textit{The Distribution of the Flora in the Alpine Zone},
  New Phytologist, vol.~11, n\textsuperscript{o}~2, pp.~37--50, 1912.

\bibitem{pagerank1998}
  L.~Page, S.~Brin, R.~Motwani et T.~Winograd,
  \textit{The PageRank Citation Ranking: Bringing Order to the Web},
  Rapport technique, Stanford University, 1998.

\bibitem{brandes2001}
  U.~Brandes,
  \textit{A Faster Algorithm for Betweenness Centrality},
  Journal of Mathematical Sociology, vol.~25, n\textsuperscript{o}~2,
  pp.~163--177, 2001.

\bibitem{tfidf}
  G.~Salton et C.~Buckley,
  \textit{Term-weighting Approaches in Automatic Text Retrieval},
  Information Processing \& Management, vol.~24, n\textsuperscript{o}~5,
  pp.~513--523, 1988.

\bibitem{gutenberg}
  Projet Gutenberg,
  \textit{Free eBooks},
  \url{https://www.gutenberg.org}, consulté en novembre 2025.

\bibitem{fastapi}
  S.~Ramírez,
  \textit{FastAPI Documentation},
  \url{https://fastapi.tiangolo.com}, consulté en novembre 2025.

\end{thebibliography}

\end{document}
